\chapter{نقشهٔ راه اجرایی}
\label{chap:roadmap}
%=====================================================================

\section{فازبندی پیشنهادی}

اجرای طرح در سه فاز اصلی و به‌صورت هم‌پوشان پیشنهاد می‌شود:

\begin{itemize}
\item \textbf{فاز \lr{1} -- امکان‌سنجی و طراحی دقیق (ماه $1$ تا $2$):} نقشه‌برداری فیبر موجود، ارزیابی تلفات واقعی مسیرهای پیشنهادی، انتخاب نهایی گره‌ها و تدوین سند الزامات امنیت فیزیکی.
\item \textbf{فاز \lr{2} -- پیوند آزمایشی دوگانه (ماه $3$ تا $5$):} استقرار نخستین پیوند \lr{QKD} میان دو گرهٔ اولویت‌دار (مثلاً بانک مرکزی و مرکز دادهٔ شتاب/شاپرک) به‌عنوان اثبات مفهوم.
\item \textbf{فاز \lr{3} -- تکمیل حلقهٔ کلان‌شهری (ماه $5$ تا $9$):} افزودن سایر گره‌ها و تکمیل توپولوژی حلقه‌ای.
\item \textbf{فاز \lr{4} -- یکپارچه‌سازی بانکی (ماه $7$ تا $11$):} اتصال \lr{KMS} به سامانه‌های پیام‌رسانی مالی موجود و آزمون تراکنش‌های واقعی در محیط کنترل‌شده.
\item \textbf{فاز \lr{5} -- بهره‌برداری و ممیزی (ماه $10$ تا $12$ و ادامه):} انتقال به بهره‌برداری کامل، ممیزی امنیتی مستقل و برنامه‌ریزی برای توسعهٔ آتی.
\end{itemize}

به‌موازات این فازها، پیاده‌سازی \textbf{لایهٔ \lr{PQC}} در نرم‌افزار (ماه $2$ تا $6$) به‌صورت مستقل از زیرساخت فیزیکی \lr{QKD} پیش می‌رود تا کوتاه‌ترین مسیر ممکن برای ارتقای امنیت کل زیرساخت بانکی کشور (حتی فراتر از محدودهٔ جغرافیایی طرح پایلوت) فراهم شود.

\begin{figure}[htbp]
\centering
\begin{ganttchart}[
  hgrid,vgrid,
  bar/.append style={fill=blue!45},
  bar height=0.6,
  title height=1,
  x unit=0.75cm,
  y unit chart=0.7cm,
]{1}{12}
\gantttitle{بازهٔ زمانی (ماه)}{12} \\
\gantttitlelist{1,...,12}{1} \\
\ganttbar{فاز \lr{1}: امکان‌سنجی}{1}{2} \\
\ganttbar{فاز \lr{2}: پیوند آزمایشی}{3}{5} \\
\ganttbar{فاز \lr{3}: تکمیل حلقه}{5}{9} \\
\ganttbar{فاز \lr{4}: یکپارچه‌سازی بانکی}{7}{11} \\
\ganttbar{لایهٔ \lr{PQC} نرم‌افزاری}{2}{6} \\
\ganttbar{فاز \lr{5}: بهره‌برداری و ممیزی}{10}{12}
\end{ganttchart}
\caption{نقشهٔ راه اجرایی فاز پایلوت (افق $12$ ماهه)}
\label{fig:gantt}
\end{figure}

\section{شاخص‌های کلیدی موفقیت \lr{(KPI)}}

\begin{itemize}
\item برقراری پایدار پیوند \lr{QKD} با نرخ خطای بیت کوانتومی کمتر از آستانهٔ امنیتی (رابطهٔ \ref{eq:qber}) به‌مدت حداقل $90$ روز متوالی؛
\item یکپارچه‌سازی موفق \lr{PQC (ML-KEM/ML-DSA)} در حداقل یک سامانهٔ پیام‌رسانی مالی داخلی؛
\item عبور موفق از یک دور آزمون نفوذ مستقل بدون یافتن آسیب‌پذیری بحرانی؛
\item تربیت حداقل $10$ متخصص داخلی در حوزهٔ طراحی، بهره‌برداری و ممیزی امنیت \lr{QKD/PQC}.
\end{itemize}

%=====================================================================
